% content.tex — Main body of the term paper
% Sections 1–7 + Appendix

\section{Introduction}

\subsection{Problem and Motivation}

Legal document-grounded question answering requires systems that can extract precise factual details from authoritative texts while producing well-reasoned explanations when a direct lookup is insufficient. In practical deployments on consumer hardware, where GPU memory is limited to a single 8\,GB card, scaling to larger language models is not an option. Instead, practitioners face a concrete engineering choice: invest in a strong retrieval-augmented generation (RAG) pipeline, apply parameter-efficient fine-tuning to the generator, or combine both.

This study investigates whether parametric adaptation retains practical value when the baseline is not a bare generator but a carefully engineered document-grounded RAG system. The question is not whether RAG helps---it demonstrably does---but whether \emph{additionally} adapting the generator's parameters yields measurable gains once the retrieval backbone is already strong. The experimental setting is deliberately constrained to consumer hardware, where retrieval engineering and parameter-efficient methods represent the most feasible adaptation strategies.

\subsection{Research Questions and Scope}

To isolate the effect of parametric adaptation from confounding factors, the study adopts a deliberately narrow scope: one compact legal benchmark, one frozen language model backbone, one hardware configuration, and one fixed retrieval stack. Within this controlled setting, two research questions guide the investigation:

\textbf{RQ1.} Does parametric adaptation yield measurable gains over a strong RAG baseline on a compact legal benchmark under consumer-hardware constraints, and how do RAFT-style supervised adaptation and supervision-free CLM continued pretraining differ as retrieval-conditioned generators?

\textbf{RQ2.} How far can pure parametric systems reach without retrieval on this benchmark, and does retrieval remain indispensable?

The benchmark comprises 8 legal documents from the Dubai International Financial Centre (DIFC) regulatory corpus and 200 human-authored QA pairs spanning six answer types. All systems share the same frozen 50-question evaluation set, while supervised training uses the remaining 150 questions. This narrow scope is methodologically intentional: it permits a controlled comparison of training signals under identical infrastructure.

\subsection{Contributions}

The paper makes three contributions:

\begin{enumerate}
    \item A controlled comparison between RAFT-style supervised adaptation and CLM continued pretraining on top of an identical RAG backbone, isolating training signal as the sole variable.
    \item A quantification of the limits of pure parametric memory by contrasting retrieval-aware systems against no-retrieval controls.
    \item A post-hoc adapter-merge result suggesting partial complementarity between the two training signals, reported as an exploratory finding.
\end{enumerate}

\subsection{Structure of the Paper}

Section~2 provides background on RAG, parameter-efficient adaptation, and the two training paradigms under comparison. Section~3 describes the benchmark corpus, hardware constraints, and the fixed retrieval backbone. Section~4 defines the compared systems and the evaluation protocol. Section~5 presents the experimental results, including aggregate comparisons, per-type analyses, and a single-document versus multi-document breakdown. Section~6 discusses the findings in light of the research questions, analyzes common error patterns, and acknowledges the study's limitations. Section~7 concludes with practical takeaways and directions for future work.

\newpage
\section{Background and Related Work}

\subsection{RAG as Nonparametric Memory}

Retrieval-augmented generation (RAG) delegates document knowledge to an external retrieval pipeline rather than requiring the generator to memorize corpus content in its parameters~\cite{lewis2020rag}. At inference time, a query is used to retrieve relevant passages from an indexed corpus, and the retrieved evidence is prepended to the generator's input context. The generator then conditions its output on both the query and the retrieved text.

For legal QA, this externalization is particularly attractive because answers must be traceable to specific regulatory provisions. A RAG system can, in principle, cite the pages from which an answer was derived, providing a form of evidence grounding that purely parametric systems lack. On consumer hardware, where model size is constrained, retrieval also compensates for the limited world knowledge that a compact model can encode in its parameters.

\subsection{Parameter-Efficient Adaptation on Consumer Hardware}

Full fine-tuning of even a 2-billion-parameter model requires storing optimizer states and gradients for all parameters, which exceeds the memory budget of an 8\,GB consumer GPU. Low-Rank Adaptation (LoRA) addresses this by freezing the pretrained weights and injecting small trainable rank-decomposition matrices into selected attention layers~\cite{hu2022lora}. QLoRA extends this approach by quantizing the frozen backbone to 4-bit NormalFloat (NF4) precision, reducing memory consumption further while preserving adaptation quality~\cite{dettmers2023qlora}. A comprehensive survey of PEFT methods, including LoRA variants and their trade-offs, is provided by Han et~al.~\cite{han2024peft}. In this study, QLoRA enables controlled adaptation of the Gemma-2-2b-it backbone within a consumer-hardware budget. Parameter-efficient adaptation constitutes part of the experimental constraint rather than an incidental implementation detail: the same QLoRA configuration is shared across both adapted systems to prevent PEFT architecture from confounding the training-signal comparison.

\subsection{RAFT-style Adaptation vs.\ CLM Continued Pretraining}

The central experimental axis of this study contrasts two training signals applied to the same PEFT architecture.

\textbf{RAFT-style supervised adaptation.} Inspired by Retrieval-Augmented Fine-Tuning~\cite{zhang2024raft}, the adapter is trained on question--answer pairs where the input includes retrieved evidence chunks (both gold and distractor passages). This directly optimizes answer generation from evidence-rich contexts, exposing the adapter to the QA task distribution. The training signal is supervised: labels are the reference answers.

\textbf{CLM continued pretraining.} The adapter is trained on the raw corpus text using a standard causal language modeling (CLM) objective---next-token prediction on all tokens. No QA labels or task-specific formatting are used. The adapter is exposed to the corpus distribution without any task-specific supervision, relying solely on the language modeling objective to absorb domain patterns.

These two paradigms represent fundamentally different assumptions about how parametric adaptation should interact with retrieval. RAFT-style training teaches the generator \emph{how to use} retrieved evidence; CLM pretraining teaches it \emph{what the corpus contains}.

\subsection{Research Gap and Positioning}

High-level discussions of parametric versus nonparametric knowledge injection are abundant in the literature. In the legal domain specifically, benchmarks such as LegalBench~\cite{guha2023legalbench}, LegalBench-RAG~\cite{pipitone2024legalbenchrag}, and evaluation frameworks like LRAGE~\cite{park2025lrage} have evaluated LLM capabilities for legal reasoning and retrieval-augmented legal QA. However, controlled studies that isolate the training signal while holding the retrieval infrastructure, PEFT architecture, backbone, and evaluation protocol fixed remain scarce. Most comparisons involve different model families, different retrieval stacks, or different evaluation setups, making it difficult to attribute performance differences to the training signal alone. This study fills that gap for a specific, practically motivated setting: compact legal QA on consumer hardware.

\newpage
\section{Benchmark and Experimental Setup}

\subsection{Corpus and Benchmark}

The benchmark is built on 8 PDF documents from the DIFC legal corpus, comprising statutes, regulations, and court judgments. Together, the documents span approximately 176 pages and 115,000 tokens. A pool of 200 question--answer pairs was authored by domain experts, covering six answer types: free-text explanations (53 questions), boolean lookups (48), numeric extractions (36), named entity lookups (30), multi-name lists (17), and date extractions (16). The distribution includes 26 multi-document comparative questions (13\%) and 17 unanswerable questions (8.5\%), ensuring that evaluation is not limited to simple single-document lookups.

The benchmark combines heterogeneous answer types---from boolean lookups to free-text legal explanations---rather than reducing evaluation to a single regime. This heterogeneity exposes distinct failure modes and prevents aggregate scores from masking type-specific weaknesses.

The 200 questions are split into 150 training questions and 50 evaluation questions, stratified by answer type, difficulty, and single-/multi-document status. All systems are evaluated on the identical 50-question evaluation set. Supervised training (RAFT-style adaptation) uses only the 150 training questions; CLM continued pretraining uses the raw document text and is therefore independent of the QA split.

\subsection{Hardware, Shared Backbone, and Variance Policy}

All experiments run on a single NVIDIA RTX 4060 with 8\,GB VRAM and 32\,GB system RAM. The language model backbone is Gemma-2-2b-it, an instruction-tuned model with approximately 2 billion parameters. The backbone is held constant across all systems to prevent architectural variance from confounding the training-signal comparison.

For systems that involve training (S2+R, S3+R, and their no-retrieval controls), three random seeds (42, 123, 777) are used, and results are reported as mean $\pm$ standard deviation. No cross-validation is performed: the single frozen split is shared across all evaluations, and seed-level variance captures only the stochasticity introduced by the training process.

\subsection{Fixed Retrieval Backbone}

The retrieval stack is held constant across all retrieval-aware systems (S1, S2+R, S3+R, S7). It comprises a five-stage pipeline:

\begin{enumerate}
    \item \textbf{Ingestion and hierarchical chunking.} Documents are parsed and split into five chunk families: page-level, section-level, clause-level, microchunks (300 tokens, 50-token overlap), and table blocks.
    \item \textbf{Hybrid retrieval.} Each query is embedded using Qwen3-Embedding-0.6B (384 dimensions) for dense retrieval and tokenized for BM25 sparse retrieval ($k_1 = 1.5$, $b = 0.75$). Both channels prefetch 30 candidates.
    \item \textbf{Reciprocal Rank Fusion (RRF).} Dense and sparse candidate lists are fused with equal weights and $k = 60$, producing a ranked list of 10 candidates.
    \item \textbf{Cross-encoder reranking.} The top 10 candidates are reranked using Qwen3-Reranker-0.6B, and the top 5 are retained.
    \item \textbf{Evidence compression.} A page-diverse compressor selects up to 3 chunks (at most one per physical page), and the corresponding (doc\_id, page\_number) pairs are lifted for grounding evaluation.
\end{enumerate}

Because the retrieval backbone is identical for all retrieval-aware systems, performance differences should be attributed to how the generator utilizes retrieved evidence, not to differences in evidence selection. Exact configuration parameters are listed in Appendix~A.

\newpage
\section{Compared Systems and Evaluation Protocol}

\subsection{System Inventory}

The study evaluates seven systems that serve distinct roles. Three form the headline comparison, one provides an exploratory post-hoc result, and three serve as negative controls. Table~\ref{tab:systems} summarizes their key characteristics.

\begin{table}[ht]
\centering
\caption{Compared systems and their roles.}
\label{tab:systems}
\begin{tabular}{lllll}
\toprule
\textbf{System} & \textbf{Retrieval} & \textbf{Training Signal} & \textbf{Supervision} & \textbf{Role} \\
\midrule
S1      & Yes & None              & ---          & Headline baseline \\
S2+R    & Yes & RAFT-style QA     & Supervised   & Headline \\
S3+R    & Yes & CLM on corpus     & Unsupervised & Headline \\
S7      & Yes & Merged S2+R+S3+R  & Post-hoc     & Exploratory \\
S2      & No  & RAFT-style QA     & Supervised   & Control \\
S3      & No  & CLM on corpus     & Unsupervised & Control \\
S3-leg. & No  & D2L hypernetwork  & Supervised   & Legacy control \\
\bottomrule
\end{tabular}
\end{table}

\begin{figure}[ht]
    \centering
    \includegraphics[width=0.92\textwidth]{figures/fig01_system_schematic.png}
    \caption{System overview schematic. Retrieval-aware systems route queries through a shared retrieval backbone to adapted generators; controls bypass retrieval entirely.}
    \label{fig:schematic}
\end{figure}

\textbf{S1 (Classical RAG)} serves as the nonparametric baseline: the frozen Gemma-2-2b-it generator receives retrieved evidence and produces answers without any adapter.

\textbf{S2+R (RAFT-style QLoRA + retrieval)} and \textbf{S3+R (CLM + retrieval)} are the two headline adapted systems. Both use the same retrieval stack as S1 and the same QLoRA architecture (rank 32, alpha 32, dropout 0.05, targeting q\_proj and v\_proj). They differ only in training signal.

\textbf{S7 (Post-hoc adapter merge)} linearly interpolates the S2+R and S3+R adapters with equal weights ($\alpha = 0.5$) without any additional training.

\textbf{S2} and \textbf{S3} are parametric controls that bypass retrieval at inference time. \textbf{S3-legacy (D2L)} is a historical negative control using a Doc-to-LoRA hypernetwork approach~\cite{charakorn2026doc} that required a chunk-level workaround; engineering details are in Appendix~C.

\subsection{Training Setups}

Both S2+R and S3+R employ identical QLoRA configurations applied to the same frozen backbone; the sole difference is the training objective.

\textbf{S2+R training.} The adapter is fine-tuned for 3 epochs on the 150 training questions in RAFT format. Each training example consists of the question, gold evidence chunks, and 2 distractor chunks. The target is the reference answer. Learning rate is $2 \times 10^{-4}$ with cosine decay and 3\% warmup. Maximum sequence length is 4096 tokens. Training takes approximately 20 minutes per seed.

\textbf{S3+R training.} The adapter is pretrained for 5 epochs on the concatenated corpus text ($\sim$115K tokens) using a CLM objective. Learning rate is $5 \times 10^{-5}$ with cosine decay and 10\% warmup. Maximum sequence length is limited to 512 tokens because computing CLM cross-entropy over the full Gemma vocabulary ($\sim$256K tokens) at longer sequences exceeds 8\,GB VRAM. Training takes approximately 10 minutes per seed.

\textbf{S7.} No training is performed. The adapter weight matrices are linearly interpolated: $W_{\text{merged}} = 0.5 \cdot W_{\text{S2+R}} + 0.5 \cdot W_{\text{S3+R}}$.

\subsection{Evaluation Protocol}

The evaluation protocol combines deterministic scoring for structured answer types with judge-based assessment for free-text responses.

\textbf{Composite metric.} The primary metric is $Q_{\text{main}} = 0.7 \times S_{\text{det}} + 0.3 \times S_{\text{asst}}$, weighting deterministic extraction at 0.7 and judged free-text quality at 0.3.

\textbf{Deterministic score ($S_{\text{det}}$).} For boolean, number, name, and date questions, scoring is binary exact match. For multi-name lists, the score is the Jaccard similarity between predicted and gold name sets.

\textbf{Free-text score ($S_{\text{asst}}$).} Free-text responses are evaluated by GPT-5.4-mini (OpenAI, reasoning effort = medium) against 5 binary criteria: correctness, completeness, grounding, calibration, and clarity (following the LLM-as-judge paradigm; see Pradhan et~al.~\cite{pradhan2025llmjudge}). The per-question score is the mean of the 5 criteria.

\textbf{Grounding ($G$).} For retrieval-aware systems, grounding is computed as $F_\beta$ ($\beta = 2.5$) on page-level (doc\_id, page\_number) pairs. Because the retrieval stack is fixed, grounding primarily reflects the shared evidence pipeline rather than adapter-specific behavior. The constant $G = 0.567$ across all retrieval-aware systems is a direct consequence of this design choice.

\newpage
\section{Results}

\subsection{Main Comparison}

Table~\ref{tab:results} presents the aggregate results. The headline systems are grouped at the top, followed by the exploratory post-hoc merge, and then the negative controls.

\begin{table}[ht]
\centering
\caption{Main results on the 50-question evaluation set. Trained systems report mean $\pm$ std across 3 seeds. S7 offline cost = S2+R (1206\,s) + S3+R (581\,s).}
\label{tab:results}
\small
\begin{tabular}{llccccccc}
\toprule
 & \textbf{System} & $Q_{\text{main}}$ & $S_{\text{det}}$ & $S_{\text{asst}}$ & $G$ & \textbf{Lat.\,(ms)} & \textbf{Offline\,(s)} \\
\midrule
\multirow{3}{*}{\rotatebox{90}{\scriptsize Headline}}
 & S1      & 0.643            & 0.601           & 0.739            & 0.567 & 479  & ---  \\
 & S2+R    & 0.669$\pm$0.014  & 0.648$\pm$0.015 & 0.718$\pm$0.018  & 0.567 & 492  & 1206 \\
 & S3+R    & 0.667$\pm$0.023  & 0.599$\pm$0.016 & 0.826$\pm$0.062  & 0.567 & 525  & 581  \\
\midrule
\rotatebox{90}{\scriptsize Post}
 & S7      & 0.705$\pm$0.035  & 0.679$\pm$0.048 & 0.764$\pm$0.018  & 0.567 & 527  & 1787 \\
\midrule
\multirow{3}{*}{\rotatebox{90}{\scriptsize Control}}
 & S2      & 0.263$\pm$0.005  & 0.270           & 0.246            & ---   & 257  & 88   \\
 & S3      & 0.185$\pm$0.003  & 0.135           & 0.303            & ---   & 195  & 581  \\
 & S3-leg. & 0.210            & 0.135           & 0.385            & ---   & 179  & 3932 \\
\bottomrule
\end{tabular}
\end{table}

The nonparametric baseline S1 already reaches $Q_{\text{main}} = 0.643$, establishing a high bar. Both retrieval-aware adapters improve over S1: S2+R attains $0.669 \pm 0.014$ and S3+R attains $0.667 \pm 0.023$. The improvements are moderate---$+0.026$ for S2+R and $+0.025$ for S3+R---but consistent across seeds. S7 reaches the highest observed score ($0.705 \pm 0.035$) through post-hoc adapter interpolation but warrants cautious interpretation due to its higher seed variance and post-hoc nature.

\begin{figure}[ht]
    \centering
    \includegraphics[width=0.85\textwidth]{figures/fig02_delta_bars.png}
    \caption{Improvement over S1 baseline. RAFT (S2+R) raises $S_{\text{det}}$, CLM (S3+R) raises $S_{\text{asst}}$, and the merge (S7) raises both.}
    \label{fig:delta}
\end{figure}

\subsection{Trade-off Between S2+R and S3+R}

The headline comparison does not yield a single dominant system. S2+R achieves higher deterministic extraction scores ($S_{\text{det}} = 0.648$ vs.\ $0.599$), whereas S3+R achieves substantially higher free-text answer quality ($S_{\text{asst}} = 0.826$ vs.\ $0.718$). On the aggregate $Q_{\text{main}}$, the two systems are near-tied ($\Delta = 0.002$). S3+R also incurs lower offline cost (581\,s vs.\ 1206\,s per seed).

\begin{figure}[ht]
    \centering
    \includegraphics[width=0.88\textwidth]{figures/fig03_judge_criteria.png}
    \caption{Judge criteria profile for free-text questions. S3+R's advantage is concentrated in correctness and completeness dimensions.}
    \label{fig:judge}
\end{figure}

\subsection{By Answer Type}

Type-level analysis reveals that performance differences between systems are concentrated in specific answer categories. Table~\ref{tab:pertype} presents the breakdown.

\begin{table}[ht]
\centering
\caption{Scores by answer type (headline and exploratory systems).}
\label{tab:pertype}
\small
\begin{tabular}{lcccccc}
\toprule
 & \textbf{Bool.} & \textbf{Num.} & \textbf{Name} & \textbf{Names} & \textbf{Date} & \textbf{Free-text} \\
 & (n=12) & (n=7) & (n=8) & (n=5) & (n=5) & (n=13) \\
\midrule
S1   & 0.833 & 0.714 & 0.500 & 0.450 & 0.200 & 0.739 \\
S2+R & 0.889 & 0.714 & 0.625 & 0.261 & 0.400 & 0.718 \\
S3+R & 0.833 & 0.714 & 0.583 & 0.300 & 0.200 & 0.826 \\
S7   & 0.889 & 0.810 & 0.708 & 0.224 & 0.400 & 0.764 \\
\bottomrule
\end{tabular}
\end{table}

S2+R outperforms S1 on boolean ($+0.056$), name ($+0.125$), and date ($+0.200$) types. S3+R shows its advantage primarily on free-text ($+0.087$ vs.\ S1). Date extraction and multi-name normalization remain weak spots across all systems. S7 achieves the highest score in 4 of 6 types.

\begin{figure}[ht]
    \centering
    \includegraphics[width=0.92\textwidth]{figures/fig04_per_type_heatmap.png}
    \caption{Score by answer type for headline and exploratory systems, with sample sizes.}
    \label{fig:heatmap}
\end{figure}

\subsection{Retrieval Contribution and the Limits of Pure Parametric Memory}

Removing retrieval causes severe quality collapse: $Q_{\text{main}}$ drops from 0.669 to 0.263 for RAFT (S2+R $\to$ S2, $\Delta = 0.406$) and from 0.667 to 0.185 for CLM (S3+R $\to$ S3, $\Delta = 0.482$). The D2L legacy control ($Q_{\text{main}} = 0.210$) corroborates this finding. These results indicate that retrieval remains the dominant memory mechanism in this setting.

\subsection{Single-Document vs.\ Multi-Document Difficulty}

Multi-document questions remain substantially harder. Table~\ref{tab:multidoc} presents this breakdown.

\begin{table}[ht]
\centering
\caption{$Q_{\text{main}}$ by document scope (42 single-doc, 8 multi-doc evaluation questions).}
\label{tab:multidoc}
\begin{tabular}{lccc}
\toprule
 & \textbf{Single-doc} & \textbf{Multi-doc} & $\Delta$ \\
\midrule
S1   & 0.694 & 0.338 & $-0.356$ \\
S2+R & 0.692 & 0.529 & $-0.163$ \\
S3+R & 0.719 & 0.338 & $-0.381$ \\
S7   & 0.716 & 0.621 & $-0.095$ \\
\bottomrule
\end{tabular}
\end{table}

S3+R achieves the highest single-doc score (0.719) but offers no improvement over S1 on multi-doc questions (both 0.338). S2+R's multi-doc score (0.529) represents a 56\% relative improvement over S1. S7 partially combines both advantages, showing the smallest gap ($\Delta = -0.095$).

\begin{figure}[ht]
    \centering
    \includegraphics[width=0.85\textwidth]{figures/fig05_singledoc_multidoc.png}
    \caption{Single-document vs.\ multi-document performance for headline and exploratory systems.}
    \label{fig:multidoc}
\end{figure}

\newpage
\section{Discussion and Limitations}

\subsection{Answer to RQ1}

RQ1 receives a qualified affirmative: parametric adaptation adds value on top of strong RAG, but the gain is moderate ($+0.026$ and $+0.025$ $Q_{\text{main}}$ respectively). The choice of training signal proves more consequential than the presence of an adapter per se. RAFT-style supervision improves deterministic extraction ($S_{\text{det}}$: $+0.047$) at the cost of a slight decrease in free-text quality ($S_{\text{asst}}$: $-0.021$). CLM continued pretraining improves free-text answer quality ($S_{\text{asst}}$: $+0.087$) while leaving deterministic extraction unchanged ($S_{\text{det}}$: $-0.002$).

The post-hoc merge S7 achieves the highest aggregate score (0.705), suggesting partial complementarity. Recent work on LoRA adapter composition~\cite{prabhakar2024lorasoups} has shown that carefully designed merge schemes can approach multi-task training quality without retraining; more structured alternatives such as rank-wise clustering~\cite{zhao2024loralego} suggest further room for improvement.

\subsection{Answer to RQ2}

RQ2 receives an unambiguous answer: retrieval remains indispensable. Neither supervised closed-book adaptation (S2, $Q_{\text{main}} = 0.263$) nor corpus-level CLM pretraining without retrieval (S3, $Q_{\text{main}} = 0.185$) provides a viable substitute. The D2L legacy control ($Q_{\text{main}} = 0.210$) corroborates this.

\subsection{Error Analysis}

Error overlap analysis reveals that 15 of the 50 evaluation questions are missed by all headline systems (Jaccard overlap = 0.714). Persistent failure patterns include date extraction, multi-name list normalization, and cross-document composition. Local wins are sparse: 2 questions answered only by S1, 2 only by S3+R, and 0 only by S2+R.

\subsection{Limitations}

\begin{itemize}
    \item \textbf{Compact corpus.} 8 documents ($\sim$115K tokens). Results may not generalize to larger corpora.
    \item \textbf{Small evaluation set.} 50 questions with per-type samples as small as $n=5$.
    \item \textbf{Single backbone.} All experiments use Gemma-2-2b-it.
    \item \textbf{Fixed retrieval.} The study measures evidence-conditioned generation, not retrieval quality.
    \item \textbf{Judge-based scoring.} $S_{\text{asst}}$ depends on GPT-5.4-mini rather than human evaluation.
    \item \textbf{Post-hoc S7.} Not pre-registered; should not be treated as a validated recommendation.
    \item \textbf{D2L as legacy control.} Required an engineering workaround; included only as a historical data point.
\end{itemize}

\newpage
\section{Conclusion}

This study investigated whether parametric adaptation adds measurable value on top of a strong RAG baseline for document-grounded legal QA on consumer hardware. The main findings are:

\begin{enumerate}
    \item \textbf{A strong RAG baseline is already highly effective.} S1 achieves $Q_{\text{main}} = 0.643$.
    \item \textbf{Training signal matters more than the mere presence of an adapter.} Both S2+R and S3+R improve over S1, but along different quality axes: RAFT improves deterministic extraction, CLM improves free-text quality. Aggregate gains are moderate ($+0.026$ and $+0.025$).
    \item \textbf{Retrieval remains non-substitutable.} Pure parametric controls collapse to $Q_{\text{main}} < 0.27$.
    \item \textbf{Post-hoc adapter fusion suggests complementarity.} S7 achieves $Q_{\text{main}} = 0.705$ and the strongest multi-document performance.
\end{enumerate}

Future work should target multi-document reasoning (where even the best system achieves only 0.621 $Q_{\text{main}}$), unanswerable-question calibration, and retrieval-aware adaptation strategies that explicitly model cross-document evidence composition.

% ── Appendix ──────────────────────────────────────────────────────
\newpage
\begin{appendices}

\section{Hyperparameters and Prompts}

\begin{table}[ht]
\centering
\caption{QLoRA configuration (shared across S2+R and S3+R).}
\small
\begin{tabular}{ll}
\toprule
\textbf{Parameter} & \textbf{Value} \\
\midrule
Rank / Alpha     & 32 / 32 \\
Dropout          & 0.05 \\
Target modules   & q\_proj, v\_proj \\
Quantization     & 4-bit NF4, double quantization \\
Optimizer        & Paged AdamW 8-bit \\
Scheduler        & Cosine \\
Weight decay     & 0.01 \\
\bottomrule
\end{tabular}
\end{table}

\begin{table}[ht]
\centering
\caption{Training-signal-specific parameters.}
\small
\begin{tabular}{lcc}
\toprule
\textbf{Parameter} & \textbf{S2+R (RAFT)} & \textbf{S3+R (CLM)} \\
\midrule
Learning rate    & $2 \times 10^{-4}$ & $5 \times 10^{-5}$ \\
Epochs           & 3                  & 5 \\
Warmup ratio     & 0.03               & 0.10 \\
Max seq.\ length & 4096               & 512 \\
Effective batch  & 4                  & 4 \\
\bottomrule
\end{tabular}
\end{table}

\section{D2L Engineering Note}

The Doc-to-LoRA (D2L) approach generates document-specific LoRA adapters via a hypernetwork. In this study, the released implementation imposed stricter effective memory limits than a preliminary token-based audit suggested, necessitating chunked packaging: each document was split into chunks, a separate adapter was generated for each, and the resulting adapters were merged via linear interpolation. This multi-stage process introduced an offline cost of 3932\,s ($\sim$65 minutes) and yielded $Q_{\text{main}} = 0.210$, placing it far below any retrieval-aware system. D2L was archived after the pilot evaluation and retained only as a legacy negative control.

\section{Use of Generative AI}

The following generative AI tools were used during the preparation of this work:

\begin{itemize}
    \item \textbf{Claude (Anthropic):} Experiment orchestration, code generation for evaluation and training pipelines, data analysis, and writing assistance (structure development, draft editing).
    \item \textbf{GPT-5.4-mini (OpenAI):} Used as the judge model for free-text answer evaluation ($S_{\text{asst}}$ scoring). The judge prompt is reproduced in the project repository.
\end{itemize}

All AI-generated content was reviewed, verified, and revised by the author. The author bears full responsibility for the accuracy and integrity of the final text, including all experimental results, interpretations, and conclusions.

\end{appendices}
